\documentclass{oxmathproblems}

\printanswers

\course{Understanding Analysis, 2nd Edition}
\sheettitle{Chapter 5 - Section 5.3 - The Mean Value Theorems} 

\begin{document}
\begin{questions}

  %------------------------- Problem 1 -------------------------
\miquestion Recall from Exercise 4.4.9 that a function $f: A \to \textbf{R}$ is Lipschitz on $A$ if there exists an $M > 0$ such that
\[\left|\frac{f(x)-f(y)}{x-y}\right| \leq M\]
for all $x \neq y$ in $A$.
\begin{enumerate}[label=(\alph*)]
  \item Show that if $f$ is differentiable on a closed interval $[a, b]$ and if $f'$ is continuous on $[a, b]$, then $f$ is Lipschitz on $[a, b]$.
  \item Review the definition of a contractive function in Exercise 4.3.11. If we add the assumption that $|f'(x)| < 1$ on $[a, b]$, does it follow that $f$ is
  contractive on this set?
\end{enumerate}

\begin{solution}
  \begin{enumerate}[label=(\alph*)]
    \item Continuous function preserves compactness (Theorem 4.4.1) and compact sets are bounded. Since $f'$ is continuous on $[a, b]$, it's bounded.
    Thus there exists $M > 0$ such that $|f'(c)| \leq M$ for all $c \in [a, b]$.

    Now, by the Mean Value Theorem, for any $x \neq y$ in $[a, b]$, there exists $t$ in $[a, b]$ such that
    \[\frac{f(x)-f(y)}{x-y} = f'(t)\] 
    Therefore
    \[\left|\frac{f(x)-f(y)}{x-y}\right| = |f'(t)| \leq M\]
    as required.

    \item $f$ is contractive if there exists a constant $0 \leq c < 1$ such that $|f(x)-f(y)| \leq c|x-y|$ for all $x, y \in \textbf{A}$.
    
    The answer is yes. From part (a), we have
    \[\left|f(x)-f(y)\right| = |f'(t)||x-y|\]
    And we have $|f'(t)| < 1$.
    
    But we need to choose a single constant $c$ (whereas $|f'(t)|$ depends on $x$ and $y$). From part (a), $f'$ is continuous on a compact
    set $[a, b]$ so it attains a maximum. Choose $c = \max_{x \in [a,b]}|f'(x)|$. Since $|f'(x)| < 1$ on $[a, b]$ from the problem's assumption, the contractive condition
    is satisfied.
  \end{enumerate}
\end{solution}

%------------------------- Problem 2 -------------------------
\miquestion Let $f$ be differentiable on an interval $A$. If $f'(x) \neq 0$ on $A$, show that $f$ is one-to-one on $A$. Provide an example to show that the converse statement
need not be true.
\begin{solution}
  We prove the contrapositive. Suppose $f$ is not one-to-one, then there exists $a < b$ in $A$ such that $f(a)=f(b)$. By the Mean Value Theorem, there exists $c \in A$ such that
  \[\frac{f(b)-f(a)}{b-a} = 0 = f'(c)\]
  so $f'$ cannot be non-zero on all of $A$.

  Counterexample to the converse: $x^2$ on $[0, 1]$. It's one-to-one (because we exclude negative domain) with $f'(0)=0$.

  Alternatively, a cleaner counterexample that doesn't involve one-sided derivative is $x^3$.
\end{solution}

\end{questions}
\end{document}
